\documentclass{situationist-preprint}

\title{Co-Endorsement Is Not Circular}
\msid{REL-09-b}
\msfamily{REL-09}
\msdiscipline{Philosophy of Religion}
\mstype{Companion manuscript}
\mskeywords{Euthyphro dilemma; co-endorsement; grounding; circularity; natural goodness; divine command}

\begin{document}
\maketitle

\begin{abstract}
REL-09-a resolves the Euthyphro dilemma by grounding both goodness and divine command in co-endorsement: an act-type's being endorsed by God and by the natures of the creatures it concerns. The account is charged with circularity in three forms: that the definition of co-endorsement uses a normative notion, flourishing, and so presupposes goodness; that God's endorsement must itself be grounded in goodness; and that we can identify what is co-endorsed only by first knowing what is good. We show that the grounding structure of the account is acyclic, that flourishing is a nature-relative notion distinct from the moral goodness of act-types, and that the third charge confuses the order of knowing with the order of grounding.
\end{abstract}

\section{The structure of grounds}

REL-09-a's account has five kinds of fact for an act-type $a$: God's endorsing $a$ ($e_G$), creaturely natures' endorsing $a$ ($e_N$), $a$'s being co-endorsed ($c$), $a$'s being good ($g$), and God's commanding $a$ ($m$). God's endorsement issues from God's nature (Definition~2.1), and creaturely endorsement from the natures of the kinds concerned. The grounding relations it asserts are that $e_G$ and $e_N$ ground $c$ (Assumption~2.5), and that $c$ grounds $g$ and $m$ (Assumption~2.6), with every ground of $g$ or of $m$ being $c$ or a ground of $c$ (Assumption~2.7).

\begin{proposition}\label{p:acyclic}
The grounding relations of REL-09-a form an acyclic order: no fact among $e_G$, $e_N$, $c$, $g$, $m$ grounds itself through a chain of these relations.
\end{proposition}

\begin{proof}
Assign ranks $0$ to $e_G$ and $e_N$, $1$ to $c$, and $2$ to $g$ and $m$. Each grounding relation of the account goes from a fact of lower rank to one of higher rank, and so does every chain of them; a chain from a fact to itself would go from a rank to the same rank. By Assumption~2.7, nothing further grounds $g$ or $m$ except through $c$, so no relation added by the account closes a cycle.
\end{proof}

A circular account, in the sense relevant to grounding \citep{Fine2012}, is one on which a fact is among its own grounds. Proposition~\ref{p:acyclic} shows that the account is not circular in this sense, and it is the structural fact on which Theorem~3.2 of REL-09-a depends: neither horn holds because command and goodness sit at the same rank, each grounded in $c$ and neither in the other.

\section{Three charges}

\begin{objection}[Flourishing presupposes goodness]
Creaturely endorsement is defined by what the flourishing of creatures calls for, and flourishing is faring well, a kind of goodness. Goodness is therefore among the grounds of co-endorsement, and so of itself.
\end{objection}

\begin{reply}
The goodness at issue in the dilemma is the moral goodness of act-types, the fact $g$. Flourishing is goodness of a different kind: the condition in which a living thing realizes the form of life of its kind, a fact about what is good \emph{for} a creature given its nature \citep{Foot2001,Thompson2008}. That a wolf that does not hunt with the pack is defective, or that a human being deprived of friendship fares badly, is a fact about natures, and it does not include any fact about what acts are morally good. Grounding moral goodness partly in such facts is no more circular than grounding the health of an organism in its physiology.
\end{reply}

\begin{objection}[God endorses because it is good]
A perfectly good God endorses $a$ because $a$ is good. So $g$ grounds $e_G$, and the account has a cycle after all.
\end{objection}

\begin{reply}
If $g$ grounded $e_G$, then by Assumption~2.5, Assumption~2.6 and transitivity, $g$ would ground itself, contrary to irreflexivity (REL-09-a, Assumption~2.4); the account rules the claim out, and must give an alternative explanation of God's endorsement. It does: God endorses from God's nature, which is loving, and a loving will endorses what the flourishing of the beloved calls for. This is not the first horn, since God's endorsement is not arbitrary, being determined by God's love and by what creatures are; and it is not the second, since no moral fact about $a$ stands prior to it.
\end{reply}

\begin{objection}[Epistemic circularity]
We can find out what is co-endorsed only by finding out what is good. The account is therefore idle, or circular as an explanation.
\end{objection}

\begin{reply}
The order of knowing need not follow the order of grounding; we learn of causes from their effects without supposing that the effects cause them. That we often know what is good before we know that it is co-endorsed is compatible with co-endorsement's being what makes it good. It is also what MTE-20 predicts: divine commands confirm rather than create obligations, and confirmation is an epistemic role that runs from command to goodness without any grounding relation between them.
\end{reply}

\section{Conclusion}

The grounding order of co-endorsement is acyclic, its definition uses natural goodness rather than the moral goodness it grounds, and the dependence of our knowledge of co-endorsement on our knowledge of goodness concerns evidence, not grounds. The charge of circularity does not apply.

\begin{thebibliography}{99}\small
\bibitem[Fine(2012)]{Fine2012} Fine, K. (2012). Guide to ground. In F.~Correia and B.~Schnieder (Eds.), \emph{Metaphysical Grounding: Understanding the Structure of Reality}, 37--80. Cambridge University Press.
\bibitem[Foot(2001)]{Foot2001} Foot, P. (2001). \emph{Natural Goodness}. Clarendon Press.
\bibitem[Thompson(2008)]{Thompson2008} Thompson, M. (2008). \emph{Life and Action: Elementary Structures of Practice and Practical Thought}. Harvard University Press.
\end{thebibliography}

\end{document}
